\subsection{对称性与自伴随性}\label{sec:symmetry-and-self-adjointness}\label{sec:adjoint-operators}

\subsubsection{伴随算子}\label{subsec:banach-adjoint-operator}

现在定义算子的\emph{伴随}.

\begin{Definition}
\label{def:banach-adjoint-operator}
设 $(\mathcal{D}(A),A)$ 为从 Banach 空间 $X$ 到 Banach 空间 $Y$ 的算子, 且 $\mathcal{D}(A)$ 在 $X$ 中稠密. 将 $\mathcal{D}(A^\times)$ 定义为满足以下条件的所有 $v\in Y^*$ 所组成的集合: 存在 $w\in X^*$, 使得对于所有 $u\in\mathcal{D}(A)$ 均有
\begin{equation}\label{eq:banach-adjoint-dual-pairing-definition}
v(Au)=w(u)
\end{equation}
注意, 由于 $\mathcal{D}(A)$ 稠密, $w$ 由 $v\in\mathcal{D}(A^\times)$ 和 \eqref{eq:banach-adjoint-dual-pairing-definition} 唯一确定. 因而可用下式定义从 $Y^*$ 到 $X^*$ 的算子 $(\mathcal{D}(A^\times),A^\times)$:
\begin{equation}\label{eq:banach-adjoint-action-definition}
A^\times(v):=w
\end{equation}
其中 $v\in\mathcal{D}(A^\times)$. 称 $(\mathcal{D}(A^\times),A^\times)$ 为 $(\mathcal{D}(A),A)$ 的\textbf{伴随}.
\end{Definition}

显然 $\mathcal{D}(A^\times)$ 非空, 因为 $\{0\}\in\mathcal{D}(A^\times)$. 而且由定义直接可知 $A^\times$ 为线性算子. 此外, 对有界算子可证明以下结果.

\begin{Theorem}
\label{thm:bounded-adjoint-domain-and-norm}
对于任意有界算子 $A\in\mathcal{L}(X,Y)$, 有 $\mathcal{D}(A^\times)=Y^*$, 且 $A^\times:Y^*\to X^*$ 是有界算子, 满足 $\|A^\times\|=\|A\|$.
\end{Theorem}

证明依赖于以下引理, 它是 Hahn-Banach 定理 (Theorem~\ref{thm:hahn-banach-extension}) 的直接推论.

\begin{Lemma}
\label{lem:unit-norm-norming-functional}
设 $X$ 为 Banach 空间, $\overline x$ 为 $X$ 的任意非零元素. 那么存在一个线性泛函 $l\in X^*$, 使得
\begin{equation}\label{eq:norming-functional-unit-norm-equality}
\|l\|=1\quad\text{and}\quad l(\overline x)=\|\overline x\|.
\end{equation}

\end{Lemma}

\begingroup
\postdisplaypenalty=10000
\renewcommand{\endProof}{\par\nobreak\hfill$\blacksquare$\par\vspace{0.1cm}}
\begin{Proof}
设 $M:=\{\alpha\overline x\mid\alpha\in\mathbb{R}\}$ 为 $\overline x$ 张成的子空间. 用下式定义 $M$ 上的线性泛函 $\widetilde l$:
\begin{equation}\label{eq:norming-functional-one-dimensional-definition}
\widetilde l(\alpha\overline x)=\alpha\|\overline x\|.
\end{equation}
容易看出, $\widetilde l$ 的范数为 $1$. Hahn-Banach 定理保证 $\widetilde l$ 在整个 $X$ 上有扩张 $l$, 其范数小于或等于 $1$. 由于 $l(\overline x)=\widetilde l(\overline x)=\|\overline x\|$, 可见其范数实际上等于 $1$, 引理得证.
\end{Proof}
\endgroup

现在证明 Theorem~\ref{thm:bounded-adjoint-domain-and-norm}.

\begingroup
\postdisplaypenalty=10000
\renewcommand{\endProof}{\par\nobreak\hfill$\blacksquare$\par\vspace{0.1cm}}
\begin{Proof}
对任意有界线性泛函 $v\in Y^*$, 可见
\begin{equation}\label{eq:bounded-adjoint-induced-functional}
u\mapsto v(A(u))=:w(u)
\end{equation}
是从 $X$ 到 $\mathbb{R}$ 的线性映射. 还可看出这个映射有界, 因为
\begin{equation}\label{eq:bounded-adjoint-functional-estimate}
|w(u)|=|v(A(u))|\leq\|v\|\|A(u)\|\leq\|v\|\|A\|\|u\|.
\end{equation}
因而 $v\in\mathcal{D}(A^\times)$, 且 $w=A^\times(v)$. 由 \eqref{eq:bounded-adjoint-functional-estimate} 还可得
\begin{equation}\label{eq:bounded-adjoint-image-norm-estimate}
\|A^\times(v)\|\leq\|A\|\|v\|.
\end{equation}
因而 $\|A^\times\|\leq\|A\|$. 现在, 由前面的引理, 对每个 $\overline u\in X$, 存在 $\overline v\in Y^*$, 使得 $\|\overline v\|=1$ 且 $\overline v(A(\overline u))=\|A(\overline u)\|$. 利用伴随的定义得到\[\begin{aligned}
\|A(\overline u)\|&=\overline v(A(\overline u))\\
&=(A^\times(\overline v))(\overline u)\\
&\leq\|A^\times(\overline v)\|\|\overline u\|\\
&\leq\|A^\times\|\|\overline v\|\|\overline u\|\\
&=\|A^\times\|\|\overline u\|.
\end{aligned}\]在最后一个等式中使用了 $\|\overline v\|=1$ 这一事实. 由于 $\overline v$ 任意, 就有 $\|A\|\leq\|A^\times\|$, 证明完毕.
\end{Proof}
\endgroup
